% ---------------------------------------------------------------------
\clearpage
\section{Catálogo de características}\label{ap:catalogo}

Las 72 características se calculan sobre el latido medio de cada paciente (70 muestras por derivación, pico R en la muestra 25, nivel del segmento PR restado). Su implementación está en \var{src/brugada/features.py}.

\subsection*{Características globales y de ritmo (5)}

\begin{center}
\small
\begin{tabularx}{\linewidth}{@{}p{3.6cm}Lp{1.5cm}@{}}
\toprule
\textbf{Nombre} & \textbf{Definición} & \textbf{Unidad} \\
\midrule
\var{qrs\_onset\_ms} & Inicio del QRS respecto al pico R & ms \\
\var{qrs\_offset\_ms} & Final del QRS (punto J) respecto al pico R & ms \\
\var{qrs\_dur\_ms} & Duración del QRS & ms \\
\var{hr\_bpm} & Frecuencia cardíaca media del registro & lpm \\
\var{rr\_cv} & Coeficiente de variación de los intervalos RR & --- \\
\bottomrule
\end{tabularx}
\end{center}

\subsection*{Características por derivación (16 $\times$ 4 = 64)}

Se calculan en V1, V2, V3 y II, con el nombre de la derivación como prefijo (por ejemplo, \var{V1\_st\_slope}).

\begin{center}
\small
\begin{tabularx}{\linewidth}{@{}p{3.6cm}Lp{1.5cm}@{}}
\toprule
\textbf{Sufijo} & \textbf{Definición} & \textbf{Unidad} \\
\midrule
\var{st\_j} & Amplitud en el punto J (final del QRS) & mV \\
\var{st\_j40} & Amplitud a J+40 ms & mV \\
\var{st\_j80} & Amplitud a J+80 ms & mV \\
\var{st\_slope} & Pendiente del ST entre J y J+80 ms & mV/ms \\
\var{st\_area} & Área del ST entre J y J+80 ms & mV·ms \\
\var{st\_max} & Máximo del ST entre J y J+120 ms & mV \\
\var{t\_amp} & Amplitud con signo del extremo de la onda T (de J+100 ms al final) & mV \\
\var{t\_pos} & Posición del extremo de la onda T respecto al pico R & ms \\
\var{t\_neg} & Mínimo de la onda T & mV \\
\var{t\_negative} & Indicador de T negativa: mínimo $< -0{,}1$ mV y mayor en valor absoluto que la mitad del máximo & 0/1 \\
\var{r\_amp} & Máximo del QRS (onda R) & mV \\
\var{s\_amp} & Mínimo del QRS (onda S) & mV \\
\var{qrs\_area} & Área absoluta del QRS & mV·ms \\
\var{rprime\_amp} & Máximo del QRS posterior a la onda S (onda r') & mV \\
\var{qrs\_end\_slope} & Pendiente desde el extremo del QRS hasta el punto J (forma convexa) & mV/ms \\
\var{st\_r\_ratio} & Cociente entre la amplitud a J+40 ms y la máxima amplitud absoluta del QRS & --- \\
\bottomrule
\end{tabularx}
\end{center}

\subsection*{Agregados de V1 y V2 (3)}

\begin{center}
\small
\begin{tabularx}{\linewidth}{@{}p{3.6cm}Lp{1.5cm}@{}}
\toprule
\textbf{Nombre} & \textbf{Definición} & \textbf{Unidad} \\
\midrule
\var{v12\_st\_j40\_max} & Máximo de \var{st\_j40} entre V1 y V2 & mV \\
\var{v12\_t\_neg\_min} & Mínimo de \var{t\_neg} entre V1 y V2 & mV \\
\var{v12\_st\_r\_ratio\_max} & Máximo de \var{st\_r\_ratio} entre V1 y V2 & --- \\
\bottomrule
\end{tabularx}
\end{center}

Además se calculan dos indicadores de calidad (número de latidos conservados y correlación media con la plantilla), que se usan en el control de calidad pero no como variables predictivas.
